\section{Introduction}
\label{sec:intro}
\vspace{-1mm}

Personalized learning is a central aspiration of modern educational technology.
Students differ in prior knowledge, learning speed, and retention capacity, yet most practice platforms assign questions uniformly or via fixed curricula that ignore individual knowledge trajectories~\citep{vanlehn2011relative,doroudi2019where}.
A deeper challenge is that human memory is inherently \emph{non-stationary}: knowledge of unpracticed topics decays over time following well-established forgetting curves~\citep{ebbinghaus1885memory,murre2015replication}.
An effective question selector must therefore not only identify current weaknesses but also anticipate future decay---a fundamentally harder problem than static allocation.
The question \emph{``which topic should we quiz next?''} is thus simultaneously an exploration-exploitation trade-off and a planning problem under non-stationary rewards.

Existing approaches to adaptive question selection draw on multi-armed bandits (MAB)~\citep{auer2002finite,clement2015multi}, Thompson Sampling~\citep{thompson1933likelihood,rafferty2019statistical}, and spaced repetition heuristics such as Leitner boxes~\citep{leitner1972so}.
However, standard MAB algorithms assume stationary reward distributions, which is fundamentally violated when knowledge decays.
Sliding-window~\citep{garivier2011upper} and discounted UCB variants partially address non-stationarity but do not model the \emph{structure} of forgetting.
Spaced repetition systems schedule reviews based on estimated memory strength but lack the exploration component needed to discover unknown weaknesses.
Knowledge tracing models~\citep{corbett1994knowledge,piech2015deep} provide sophisticated student models but typically pair them with greedy or heuristic selection policies without formal exploration guarantees.

\begin{wrapfigure}{r}{0.42\textwidth}
  \centering
  \vspace{-4mm}
  \vspace{-\intextsep}
  \includegraphics[width=0.40\textwidth]{fig_teaser.png}
  \caption{Knowledge trajectories at $K\!=\!6$, $\lambda\!=\!0.01$. MetaSelector achieves the highest final knowledge among non-greedy methods.}
  \label{fig:teaser}
  \vspace{-4mm}
\end{wrapfigure}

We present \textbf{\ours}, a comprehensive framework and benchmark for adaptive question selection under forgetting.
We formalize the problem as a \emph{restless multi-armed bandit} (RMAB)~\citep{whittle1988restless}, where each knowledge topic is an arm whose state evolves even when not selected (due to forgetting).
Within this formulation, we develop four novel algorithms: \textbf{F-UCB}, which integrates decay-aware urgency into the UCB exploration bonus; \textbf{BKT-Bandit}, which maintains Bayesian posteriors over knowledge with integrated forgetting dynamics; \textbf{Advantage Index}, which computes restless bandit indices via MDP value differences; and \textbf{MetaSelector}, a UCB-over-experts meta-algorithm that selects among base selectors online, enabling principled regime interpolation.

Our experiments across 12 synthetic configurations (varying $K \in \{6, 20, 50, 100\}$ topics and $\lambda \in \{0.005, 0.01, 0.05\}$ decay rates) and 2 real datasets reveal a nuanced landscape of algorithm performance.
No single algorithm dominates all regimes: BKT-Bandit achieves top-3 in 10/12 configurations with average rank 2.50, while F-UCB delivers the largest gains ($+42\%$ over random) at high decay.
Most strikingly, we discover a \emph{Greedy-Min catastrophe}: a myopic selector with perfect knowledge of the current state that always quizzes the weakest topic finishes last in 9/12 configurations, losing to random selection in 10/12, because greedy exploitation of known states systematically neglects decaying arms.

In summary, our contributions are:
\begin{enumerate}[leftmargin=*,itemsep=1pt,topsep=2pt]
  \item \textbf{RMAB formulation with theoretical analysis.} We formalize educational quizzing as a restless bandit with Ebbinghaus forgetting, prove a precise catastrophe condition for myopic policies (Theorem~\ref{thm:catastrophe}), and establish convergence guarantees for our proposed algorithms (Theorems~\ref{thm:bkt}--\ref{thm:fucb}).
  \item \textbf{Four novel algorithms.} We introduce F-UCB, BKT-Bandit, Advantage Index, and MetaSelector, each targeting different aspects of the forgetting-aware selection problem.
  \item \textbf{Greedy-Min catastrophe.} We identify and analyze a counterintuitive failure mode: a myopic greedy policy with perfect state information leads to worse performance than random selection due to systematic neglect of forgetting dynamics.
  \item \textbf{Comprehensive benchmark.} We evaluate 16 algorithms across 12 synthetic configurations and 2 real datasets (Duolingo: 500 students, 127K interactions; ASSISTments: 500 students, 425K interactions), providing the most extensive comparison of bandit-based question selectors to date.
\end{enumerate}
